\documentclass[11pt]{letter}
\usepackage[margin=1in]{geometry}
\signature{Rogerio Jorge\\Department of Physics, University of Wisconsin--Madison}
\address{Department of Physics\\University of Wisconsin--Madison\\Madison, Wisconsin 53706, USA\\rogerio.jorge@wisc.edu}
\begin{document}
\begin{letter}{The Editors\\Physical Review E}
\opening{Dear Editors,}
Please consider the manuscript ``Resonant conversion of dark-photon fields in a plasma with a noise-free continuum kinetic method: bounds on nonlinear suppression and the role of collisions'' for publication in Physical Review E.

Particle-in-cell simulations by Hook, Huang and Shalaby (Phys.\ Rev.\ Lett.\ 137, 071004) report that nonlinear plasma dynamics saturates the resonant conversion of dark-photon dark matter into plasma oscillations. We examine the same question with a noise-free Hermite--Fourier Vlasov solver for the Maxwell--Proca--Vlasov system in one spatial and three velocity dimensions. The manuscript reports the method, its verification and its main limit (a truncation instability of the asymmetrically weighted Hermite basis that bounds the certified time windows), and two studies: a finite dark reservoir and finite coherence time of the drive, and the role of collisions.

Within the certified windows the energy transfer follows the linear resonant law, and with collisions the absorption follows the linear Drude law. At the one drive matching a published PIC run we do not reproduce the reported saturation. The model is nonrelativistic, and the coherent quiver speed it implies is large, so the result is a bound in a restricted model and not a resolution of the cosmological question. Failed and amended gates are reported.

All code, records and the scripts that generate every figure, table and quoted number are in the public repository \texttt{uwplasma/Dark-SPECTRAX}. The work has not been published or submitted elsewhere.

\closing{Sincerely,}
\end{letter}
\end{document}
